Galileo's principle of relativity imposes the following constraints on the right-hand side of Newton's equation in an inertial coordinate system:

\begin{enumerate}
\item \textbf{Time homogeneity:} Invariance under time translations implies the right-hand side does not depend explicitly on time:
$$\ddot{x} = \Phi(x, \dot{x}).$$

\item \textbf{Space homogeneity:} Invariance under spatial translations implies the right-hand side depends only on relative coordinates:
$$\ddot{x}_i = f_i(\{x_j - x_k\}).$$

\item \textbf{Galilean boost invariance:} Invariance under passage to a uniformly moving coordinate system implies the right-hand side depends only on relative velocities:
$$\ddot{x}_i = f_i(\{x_j - x_k, \dot{x}_j - \dot{x}_k\}), \quad i, j, k = 1, \ldots, n.$$

\item \textbf{Space isotropy:} Invariance under rotations implies:
$$\mathbf{F}(G\mathbf{x}, G\dot{\mathbf{x}}) = G\mathbf{F}(\mathbf{x}, \dot{\mathbf{x}}),$$
where $G$ is any orthogonal transformation.
\end{enumerate}
